\documentclass[../../main.tex]{subfiles}
\begin{document}
\section{密度}\label{section:ch16-densities}

形式的积分只在定向流形上有意义。为了能够在不可定向流形上积分，我们需要一个稍微更一般的对象。设 $V$ 是一个 $n$ 维向量空间。$V$ 上的一个\emph{密度}是一个函数 $\mu:V^n\to\mathbb{R}$，满足对每个线性映射 $A:V\to V$ 有
\begin{align*}
\mu(Av_1,\dots,A,v_n)=|\det A|\,\mu(v_1,\dots,v_n).
\end{align*}
$V$ 上所有密度构成的集合记为 $\Omega(V)$；它自身是一个一维向量空间。

\begin{definition}[流形上的密度]\index{密度}
光滑流形 $M$ 上的一个\emph{密度}是指在每点 $p\in M$ 指定 $T_pM$ 上的一个密度 $\mu_p$，其指定方式是光滑的：即在任意局部坐标下它是一个光滑函数乘以标准密度 $|\mathrm{d}x^1\cdots\mathrm{d}x^n|$。
\end{definition}

密度与体积形式之间的关系如下。如果 $M$ 是定向的，则每个体积形式 $\omega$ 决定一个密度 $|\omega|$，即 $|\omega|_p(v_1,\dots,v_n)=|\omega_p(v_1,\dots,v_n)|$。反过来，一个定向和一个密度共同决定一个体积形式。

\begin{proposition}[密度的性质]\label{proposition:ch16-properties-of-densities}
设 $M$ 是一个 $n$ 维光滑流形。
\begin{enumerate}[(i)]
\item $M$ 上存在处处非零的光滑密度。
\item 任意两个处处非零密度相差一个光滑函数因子。
\item 一个密度在任意（不定向）光滑流形上都有意义，而体积形式只在定向流形上才有意义。
\end{enumerate}
\end{proposition}

\begin{proposition}\label{proposition:ch16-density-change-of-coords}
在局部坐标 $(x^i)$ 与 $(\widetilde x^j)$ 之间，一个密度的变换规律是
\begin{align*}
\widetilde\mu=|\det(\partial x^i/\partial\widetilde x^j)|\,\mu.
\end{align*}
这与体积形式的变换相差一个绝对值符号。
\end{proposition}

\begin{proposition}\label{proposition:ch16-density-integration}
设 $\mu$ 是光滑流形 $M$ 上的一个光滑密度。我们可以无需求助于定向，就定义 $\mu$ 在 $M$ 上的积分 $\int_M\mu$，其方式与 $n$-形式的积分完全类似：在坐标卡中用标准欧氏积分定义，并用单位分解拼接。这个积分在任意微分同胚下不变。
\end{proposition}

\begin{proposition}[Riemann 密度]\label{proposition:ch16-riemannian-density-form}
设 $(M,g)$ 是一个 Riemann 流形（不一定定向）。在每个切空间 $T_pM$ 上用 $g_p$ 的体积定义一个密度，这就给出 $M$ 上一个规范的光滑密度 $\mathrm{d}V_g$。它在局部坐标下为
\begin{align*}
\mathrm{d}V_g=\sqrt{|\det g_{ij}|}\,|\mathrm{d}x^1\cdots\mathrm{d}x^n|.
\end{align*}
当 $M$ 可定向时，这与定向 Riemann 体积形式的绝对值一致。
\end{proposition}

于是，即使在不可定向流形上，我们也可以用 Riemann 密度来定义长度、面积、体积等几何量，而无需选择定向。
\end{document}
